\begin{theorem}[Stokes 幅值三次与 Lee--Yang 三次的共同二维 Artin 表示]\label{thm:xi-terminal-zm-stokes-leyang-shared-artin-representation}
沿用定理 \ref{thm:xi-terminal-zm-b2-kappa2-rational-interconstruction} 的记号，记
\[
b:=B^2,
\qquad
u:=\kappa_*^2.
\]
再沿用结论 \ref{con:xi-terminal-zm-leyang-cubic-field-isomorphism} 的生成元 \(\alpha,\beta\)。
则有共同三次数域识别
\[
\boxed{\;
\QQ(b)=\QQ(u)=\QQ(\lambda_*)=\QQ(\alpha)=\QQ(\beta)=:K.\;}
\]
令 \(L\) 为 \(K/\QQ\) 的伽罗瓦闭包。
记 \(\rho_B\) 为由 Stokes 幅值三次
\[
162b^3+1593b^2+1998b+1184=0
\]
所定义的二维 Artin 表示；等价地，\(\rho_B\) 是 \(\Gal(L/\QQ)\cong S_3\) 的标准二维不可约表示沿自然商映射
\(\Gal(\overline\QQ/\QQ)\twoheadrightarrow \Gal(L/\QQ)\)
的拉回。
再记 \(\rho_{\mathrm{LY}}\) 为定理 \ref{thm:xi-time-part9c-leyang-cubic-artin-weight1} 中附着于 Lee--Yang 三次数域 \(K\) 的二维 Artin 表示。
则
\[
\boxed{\;\rho_B\cong \rho_{\mathrm{LY}}.\;}
\]
特别地，
\[
L(s,\rho_B)=L(s,\rho_{\mathrm{LY}}),
\qquad
\det(\rho_B)=\det(\rho_{\mathrm{LY}})=\chi_{-111},
\]
其中 \(\chi_{-111}\) 对应二次分辨子域 \(\QQ(\sqrt{-111})\)。
\end{theorem}
\begin{proof}
由定理 \ref{thm:xi-terminal-zm-b2-kappa2-rational-interconstruction}，
\[
\QQ(b)=\QQ(u)=\QQ(\lambda_*).
\]
又由定理 \ref{thm:xi-terminal-zm-stokes-t-elliptic-branch-coordinate}，\(\lambda_*\) 满足
\[
16\lambda_*^3-9\lambda_*^2+1=0.
\]
结合结论 \ref{con:xi-terminal-zm-leyang-cubic-field-isomorphism} 的
\[
\QQ(\alpha)=\QQ(\beta),
\qquad
16\alpha^3-9\alpha^2+1=0,
\]
可知 \(\QQ(\lambda_*)=\QQ(\alpha)=\QQ(\beta)\)。
遂得第一条共同域识别。

再由定理 \ref{thm:xi-terminal-zm-stokes-amplitude-square-cubic-rigidity} 与定理 \ref{thm:xi-time-part9c-leyang-cubic-artin-weight1}，
Stokes 幅值三次与 Lee--Yang 三次都给出同一非伽罗瓦三次域 \(K\)，其伽罗瓦闭包均为 \(L\)，并满足
\[
\Gal(L/\QQ)\cong S_3.
\]
而 \(S_3\) 的二维不可约复表示只有一个同构类，即标准表示。
因此由 Stokes 幅值三次得到的 \(\rho_B\) 与 Lee--Yang 三次得到的 \(\rho_{\mathrm{LY}}\) 必同构：
\[
\rho_B\cong \rho_{\mathrm{LY}}.
\]

最后，定理 \ref{thm:xi-time-part9c-leyang-cubic-artin-weight1} 已给出
\[
\det(\rho_{\mathrm{LY}})=\chi_{-111}.
\]
另一方面，由推论 \ref{cor:xi-terminal-zm-stokes-b2-quadratic-resolvent}，Stokes 幅值三次的唯一二次分辨子域同样是 \(\QQ(\sqrt{-111})\)，故
\[
\det(\rho_B)=\chi_{-111}.
\]
表示同构即推出 Artin \(L\)-函数相同，得到
\[
L(s,\rho_B)=L(s,\rho_{\mathrm{LY}}).
\]
证毕。
\end{proof}

\endinput
